% NichidaiMasterExam_R6_1st_2
\begin{tcolorbox} %%R6_1st_2
	$\fbox{2}$ $A= \begin{pmatrix}1& i& i\\ i& i& 1\end{pmatrix}$とし、$X$を複素行列とする。
\begin{enumerate}
	\item $A$の随伴行列$A^{*}$を求めよ。
	\item $A^{*}A$を計算せよ。
	\item エルミート行列の定義を述べよ。また$X^{*}X$がエルミート行列であることを示せ。
\end{enumerate}
\end{tcolorbox}

\begin{souanswer} %%R6_1st_2_ans
\begin{enumerate}
	\item 随伴行列$A^{*}$は
\begin{equation*}
	A^{*}= 
\begin{pmatrix}
	\bar{1}& \bar{i}\\
	\bar{i}& \bar{i}\\
	\bar{i}& \bar{1}
\end{pmatrix}
	=
\begin{pmatrix}
	1& -i\\
	-i& -i\\
	-i& 1
\end{pmatrix}
\end{equation*}
	\item
\begin{align*}
	A^{*}A
	&=
\begin{pmatrix}
	1& -i\\
	-i& -i\\
	-i& 1
\end{pmatrix}
\begin{pmatrix}
	1& i& i\\
	i& i& 1
\end{pmatrix}\\
	&= 
\begin{pmatrix}
	2& 1+ i& 0\\
	1- i& 2& 1- i\\
	0& 1+ i& 2
\end{pmatrix}
\end{align*}
	\item エルミート行列の定義は以下
\begin{tcolorbox}
	随伴行列と元の行列が等しい正方行列。すなわち正方行列$A$に対して$A^*= \bar{A}^\mathrm{T}= A$となる行列。
\end{tcolorbox}
	任意の複素行列$X$に対して$H= X^*X$とおく。その随伴行列を$H^*$を計算する。随伴行列の性質$(AB)^*= B^*A^*$および$(X^*)^*= X$より
\begin{equation*}
	(X^*X)^*= X^*(X^*)^*= X^*X
\end{equation*}
	$(X^*X)^*= X^*X$が成り立つため$X^*X$はエルミート行列
\end{enumerate}
\end{souanswer}